\section{The filled scalar system and its coercive identity}
\label{sy:section}

We work on one fixed strict exterior obtained in
Proposition~\ref{en:rest-preparation}.  Thus $K$ is compactly supported,
the end has smooth asymptotic symbol estimates, and, for some
$0<\beta<1$ and $c>0$,
\begin{equation}\label{sy:strict-margin}
 \mu+J(Y)-I_g(Y)\ge c r^{-3-\beta},\qquad |Y|_g\le1.
\end{equation}
Here and throughout this section the constraint densities are geometric:
$2\mu=R_g+\tau^2-|K|^2$ and
$J=\operatorname{div}(K-\tau g)$, where $\tau=\operatorname{tr}_gK$.
The function $r$ agrees with the coordinate radius sufficiently far out
and is smooth and positive elsewhere.  We set $k=2$; in particular,
every transverse block below has dimension two.

\subsection{The auxiliary filling and the unknowns}

The construction also permits a fixed sign $\sigma_i\in\{1,-1\}$ on
each boundary component $S_i$, with
$H_{S_i}+\sigma_i\operatorname{tr}_{S_i}K<0$.
For the future-trapped case of Theorem~\ref{thm:numerical}, all signs are $+1$.
Allowing the signs separately will be useful when another obstacle is
inserted.  The signs are fixed during the construction.

\begin{lemma}[Filling with controlled geometry]\label{sy:filled-setup}
For every sufficiently large integer $N$, the prepared exterior is
contained in a smooth complete, boundaryless, one-ended manifold
$\Omega_N$.  Its metric and symmetric tensor extend $g,K$ and have the
following properties.  Each $S_i$ has a fixed inner collar on whose
cross-sections $H+\sigma_i\operatorname{tr}_{\rm tan}K<0$, with the
normal directed toward the exterior.  Behind this collar there is a
fixed transition to a product cylinder, a product neck of length
$O(1+N)$, and a fixed cap.  On the product neck and cap,
$K=-\sigma_iL_0g$ for a fixed sufficiently large $L_0$.
Local coordinate bounds for the metric, its inverse, and the tensor,
of every fixed derivative order, and a positive local injectivity
radius can be chosen independently of $N$.  The enlarged compact core
has volume $O(1+N)$.  The extension of $r$ can be chosen uniformly
comparable to one there.  The closed flux forms $\alpha_1,\alpha_2$
are retained on the original exterior.
\end{lemma}

\begin{proof}
Every connected component of the boundary is an orientable closed
surface, hence is diffeomorphic to the boundary of a handlebody of
its genus.  Attach a separate such handlebody on each component.
The collar theorem gives signed collar coordinates, increasing toward
the original exterior.  Extend the metric and tensor smoothly a short
distance inward.  The strict expansion inequality persists on a
shorter collar by continuity.

On the next fixed collar, interpolate the metric to a product metric.
The metrics in this interpolation are positive definite and form a
fixed compact family; their cross-sectional mean curvatures are
therefore bounded.  Before making this interpolation, subtract
$\sigma_i Lg$ from $K$, with $L\ge0$ increasing smoothly from zero.
This changes $\sigma_i\operatorname{tr}_{\rm tan}K$ by $-2L$.
Start the change within the region where the expansion is already
strict; after $L$ has become large, it dominates the bounded mean
curvature and the bounded tensor terms in the interpolation.
One may then interpolate the remaining tensor to the pure trace
$-\sigma_iL_0g$, still retaining strictness.  Finish with an exactly
product segment.  This gives the asserted fixed transition.

Choose a smooth metric on each handlebody extending its product
boundary collar, and put $K=-\sigma_iL_0g$ on that cap.  Insert between
the transition and cap as many copies of the fixed product segment as
are needed.  A length bounded by a fixed multiple of $1+N$ is
available, with the multiple enlarged when the height barriers are
chosen.  There are finitely many fixed transition and cap pieces.
Product translation identifies all local neck charts, so all asserted
local bounds and the local injectivity radius are uniform.  Only
the cylinder volume grows with its length.  The original end is
unaltered, proving completeness and the one-end assertion.
Extending $r$ as a positive constant along the inserted pieces proves
the weight assertion.  This construction concerns the metric and
tensor only; all subsequent uses of the flux forms take place on
the original exterior, where they remain closed with their prescribed
two fluxes.
\end{proof}

Fix $0<\epsilon<1$.  Choose once and for all a smooth nonincreasing
function $\vartheta:\mathbb R\to[0,1]$, equal to one on $(-\infty,0]$
and zero on $[1,\infty)$.  Define, for $t\in\mathbb R$,
\begin{equation}\label{sy:compression-functions}
 p(t)=N\vartheta(Nt),\qquad
 l(t)=\exp\left(\int_0^t p(s)\,ds\right),\qquad
 \ell=e^{-N\epsilon},\qquad \eta=\ell^{3/2}.
\end{equation}
We have $l'=pl$, $0\le p\le N$, $l(t)=e^{Nt}$ for $t\le0$,
and $1\le l(t)\le e$ for $t\ge0$.  In particular,
\begin{equation}\label{sy:l-bounds}
 \ell\le l(t)\le e\quad\hbox{if }t\ge-\epsilon.
\end{equation}

For unknown smooth functions $f,Z$, define $t$ by
\begin{equation}\label{sy:implicit-t}
 Z=t+\frac14\log D,\qquad D=1+l(t)^2|df|_g^2.
\end{equation}
This definition is global, including for trial functions that do not
satisfy a lower bound on $t$.  Indeed, for fixed $\xi\in T_x^*\Omega_N$,
the derivative of the right side with respect to $t$ is
$1+pv/2\ge1$, where
$v=l(t)^2|\xi|^2/(1+l(t)^2|\xi|^2)$.
Its limits are $-\infty$ and $+\infty$ at the two ends of the real
line: at the negative end $l=e^{Nt}$, and at the positive end $l$ is
constant.  The implicit function theorem therefore gives a unique
smooth $t=t(x,Z,\xi)$, with no singularity at $\xi=0$.

Here $f$ is the auxiliary graph height, $e^{2t}$ the conformal multiplier,
and $Z$ the scalar used in the divergence equation.

Vectors and covectors are identified using $g$.  Our notation is
\begin{equation}\label{sy:variables}
\begin{aligned}
 w&=lD^{-1/2}\nabla f,& a_w&=|w|,& v&=a_w^2,& d&=D^{-1}=1-v,\\
 q&=2+pv,& \chi&=\frac{2d}{q},&
 A_d&=I-w\otimes w,&
 A_\chi&=I-\frac{2+p}{q}w\otimes w,\\
 U&=l\sqrt D,& u&=e^tU,& h&=\eta f,&
 H^f&=lD^{-1/2}\nabla^2 f,\\
 C&=K+H^f,& F&=\operatorname{tr}_{A_\chi}C,&
 V&=u A_\chi(2\nabla Z+K(w,\cdot)),&
 \widehat g&=e^{2t}(g+l^2df^2).
\end{aligned}
\end{equation}
Here $\operatorname{tr}_A C=A^{ij}C_{ij}$.  Both $A_d$ and
$A_\chi$ have eigenvalue one perpendicular to $w$, and their axial
eigenvalues are $d$ and $\chi$, respectively.  Thus
$0<\chi\le d\le1$, and
\begin{equation}\label{sy:axis-comparison}
 \frac{2}{2+N}A_d\le A_\chi\le A_d.
\end{equation}
The symbols $q,d,\chi$ in this section are auxiliary scalar
coefficients; $Q$ continues to denote the total charge.

At a point with $df\ne0$, let $e=\nabla f/|\nabla f|$.  Write
\begin{equation}\label{sy:blocks}
 C=\begin{pmatrix}T&M\\M^{\mathsf t}&j\end{pmatrix},\qquad
 dt=(y,x),\qquad w=a_we,
 \qquad T^0=T-\tfrac12(\operatorname{tr}T)I_{e^\perp}.
\end{equation}
The first block acts on the two-dimensional space $e^\perp$.
Consequently $F=\operatorname{tr}T+\chi j$.
At $df=0$, choose any unit vector $e$ to write these block formulas.
All scalar and tensor expressions defined invariantly below are
independent of that choice.

\subsection{The exact scalar-curvature identity}

For symmetric tensors $A,B$, put
\[
 \mathsf P_A=A-(\operatorname{tr}A)g,\qquad
 \mathsf B(A,B)=\langle A,B\rangle
                   -(\operatorname{tr}A)(\operatorname{tr}B).
\]
Let
\begin{equation}\label{sy:stationary-tensors}
 b=-H^f-p(dt\otimes w+w\otimes dt),\qquad
 \mathsf Q=K-b=C+p(dt\otimes w+w\otimes dt),\qquad
 s_p=p+\tfrac12.
\end{equation}
Define the smooth scalar
\begin{equation}\label{sy:invariant-T}
 \mathcal T=\tfrac12\mathsf B(\mathsf Q,\mathsf Q)
 +\mathsf P_{\mathsf Q}(w,\nabla t)
 +2s_p|dt|_{A_d}^2
 +F\bigl(\operatorname{tr}C+2s_p w(t)\bigr).
\end{equation}
This definition, rather than a choice of axis, fixes $\mathcal T$
also at zero slope.

\begin{proposition}[Exact filled scalar identity]\label{sy:scalar-identity}
For every pair of smooth functions $f,Z$ on the filled manifold,
\begin{equation}\label{sy:scalar-equation}
 \frac12e^{2t}R_{\widehat g}
 =\mu+J(w)+\mathcal T+w(F)-F\tau
       -u^{-1}\operatorname{div}_g V.
\end{equation}
The quadratic expression in Equation~\ref{sy:invariant-T} has the
following exact block expansion:
\begin{equation}\label{sy:expanded-T}
\begin{aligned}
 \mathcal T={}&\tfrac12|T^0|^2+|M+s_pa_wy|^2
       -\tfrac14v|y|^2+2s_p(|y|^2+dx^2)\\
 &-\tfrac14(F-\chi j)^2+\chi j^2-\chi Fj+F^2
       +2s_p\chi ja_wx.
\end{aligned}
\end{equation}
These identities require neither a constraint inequality nor a
field extension to the filling.
\end{proposition}

\begin{proof}
First obtain the scalar identity before conformal multiplication.
On $\Omega_N\times\mathbb R$ consider the stationary Lorentzian metric
\[
 \mathbf G=-l^2(dz-df)^2+\bar g,
 \qquad \bar g=g+l^2df^2.
\]
Written in the coordinates $(x,z)$, it has spatial metric $g$, lapse
$U=l\sqrt D$, and shift $\gamma=l^2\nabla f=Uw$.
Its future unit normal to the constant-$z$ slices is
$\mathbf n=U^{-1}(\partial_z-\gamma)=U^{-1}\partial_z-w$.
With the convention for $K$ in Definition~\ref{def:data}, stationarity
gives
\begin{equation}\label{sy:shift-calculation}
 b=-U^{-1}\operatorname{sym}\nabla\gamma
   =-H^f-p(dt\otimes w+w\otimes dt),\qquad
 \operatorname{div}\gamma=-U H_b,
 \quad H_b=\operatorname{tr}b.
\end{equation}
Here $\operatorname{sym}$ means the average of a two-tensor and its
transpose.  The contracted Gauss equation and normal variation of the
second fundamental form give, with this sign convention,
\[
 \mathbf R=R_g+H_b^2-|b|^2-2\mathbf{Ric}(\mathbf n,\mathbf n),
 \qquad
 \mathbf{Ric}(\mathbf n,\mathbf n)
 =-\mathbf n(H_b)-|b|^2+U^{-1}\Delta_g U.
\]
Since $H_b$ is independent of $z$, $\mathbf n(H_b)=-w(H_b)$.
Substituting and using $\operatorname{div}(Uw)=-UH_b$ yields
\begin{equation}\label{sy:stationary-scalar}
 \mathbf R
 =R_g+\mathsf B(b,b)
       -2U^{-1}\operatorname{div}(\nabla U+H_bUw).
\end{equation}
In the coordinates $z-f(x)$ the same metric is the warped static
metric $-l^2d(z-f)^2+\bar g$.  Its scalar curvature is
$\mathbf R=R_{\bar g}-2l^{-1}\Delta_{\bar g}l$.
The inverse and determinant of $\bar g$ are $A_d$ and $D\det g$,
respectively.  Hence
\begin{equation}\label{sy:warp-divergence}
 l^{-1}\Delta_{\bar g}l
 =U^{-1}\operatorname{div}\bigl((U/l)A_d\nabla l\bigr),
 \qquad
 \nabla U-(U/l)A_d\nabla l=-Ub(w,\cdot).
\end{equation}
For the second identity, divide by $U$ and use
$d\log U=p(1+v)dt+H^f(w,\cdot)$ and $d\log l=pdt$;
the resulting difference is
$H^f(w,\cdot)+p\{vdt+w(t)w\}=-b(w,\cdot)$.
Equations~\ref{sy:stationary-scalar} and
\ref{sy:warp-divergence} now give
\begin{equation}\label{sy:bar-scalar-b}
 \tfrac12R_{\bar g}
 =\tfrac12R_g+\tfrac12\mathsf B(b,b)
       +U^{-1}\operatorname{div}(U\mathsf P_b w).
\end{equation}

For clarity, the insertion of the original constraint tensor is an
exact cancellation.  Namely,
\[
 U^{-1}\operatorname{div}(U\mathsf P_Kw)
 =J(w)+\mathsf P_K:U^{-1}\nabla(Uw)
 =J(w)-\mathsf B(K,b),
\]
where symmetry of $\mathsf P_K$ permits use of
Equation~\ref{sy:shift-calculation}.  Moreover,
$\mu=R_g/2-\mathsf B(K,K)/2$ and
$\mathsf Q=K-b$.  Expanding $\mathsf B(\mathsf Q,\mathsf Q)$
in Equation~\ref{sy:bar-scalar-b} therefore proves
\begin{equation}\label{sy:bar-scalar-Q}
 \tfrac12R_{\bar g}
 =\mu+J(w)+\tfrac12\mathsf B(\mathsf Q,\mathsf Q)
        -U^{-1}\operatorname{div}(U\mathsf P_{\mathsf Q}w).
\end{equation}

In dimension three the conformal scalar law is
\[
 \tfrac12e^{2t}R_{\widehat g}
 =\tfrac12R_{\bar g}-2\Delta_{\bar g}t-|dt|_{A_d}^2.
\]
The same determinant calculation as above gives
$\Delta_{\bar g}t
=U^{-1}\operatorname{div}(UA_d\nabla t)-p|dt|_{A_d}^2$.
Set $W=\mathsf P_{\mathsf Q}w+2A_d\nabla t$.  Since $u=e^tU$,
\[
 -U^{-1}\operatorname{div}(UW)
 =-u^{-1}\operatorname{div}(uW)+W(t).
\]
It follows that
\begin{equation}\label{sy:before-flux-replacement}
\begin{aligned}
 \tfrac12e^{2t}R_{\widehat g}
 ={}&\mu+J(w)+\tfrac12\mathsf B(\mathsf Q,\mathsf Q)
   +\mathsf P_{\mathsf Q}(w,\nabla t)\\
 &+2s_p|dt|_{A_d}^2-u^{-1}\operatorname{div}(uW).
\end{aligned}
\end{equation}

Differentiating Equation~\ref{sy:implicit-t} gives the useful exact
one-form identity
\begin{equation}\label{sy:differential}
 2\,dZ=(2+pv)dt+H^f(w,\cdot).
\end{equation}
In the block notation, $\operatorname{tr}T=F-\chi j$, and
$\mathsf P_{\mathsf Q}w$ has transverse component
$a_wM+pv y$ and axial component $-a_w(F-\chi j)$.
Thus both sides of the first identity below have transverse component
$u(qy+a_wM)$ and axial component
$u(2dx+\chi a_wj-a_wF)$:
\begin{equation}\label{sy:flux-identities}
 uW=V-uFw,\qquad
 u^{-1}\operatorname{div}(uw)
 =F+(1-\chi)j-\tau+2s_pa_wx.
\end{equation}
For the second identity, use
$u^{-1}\operatorname{div}(uw)=-H_b+w(t)
=\operatorname{tr}H^f+(2p+1)w(t)$.
Replacing the flux in Equation~\ref{sy:before-flux-replacement}
and differentiating $uFw$ gives
Equation~\ref{sy:scalar-equation} with
Equation~\ref{sy:invariant-T}.

To check the full expansion, put $z_T=\operatorname{tr}T=F-\chi j$.
The transverse, mixed, and axial blocks of $\mathsf Q$ are
$T$, $M+pa_wy$, and $j+2pa_wx$.  Therefore
\[
 \tfrac12\mathsf B(\mathsf Q,\mathsf Q)
 =\tfrac12|T^0|^2-\tfrac14z_T^2
   +|M+pa_wy|^2-z_Tj-2pa_wxz_T,
\]
and
\[
 \mathsf P_{\mathsf Q}(w,\nabla t)
 =a_w\langle M,y\rangle+pv|y|^2-a_wxz_T.
\]
Combining the terms in $M,y$ completes the square
$|M+s_pa_wy|^2-v|y|^2/4$.
The remaining linear terms in $x$ equal
$2s_pa_wx(F-z_T)=2s_p\chi ja_wx$;
the remaining terms in $j,F$ equal
$-z_Tj+F(F+(1-\chi)j)=\chi j^2-\chi Fj+F^2$.
These are exactly all the terms in Equation~\ref{sy:expanded-T}.
\end{proof}

\subsection{Polynomial coercivity and the stationary-point comparison}

We use $\mathcal P_N$ for a positive bound of the form
$C(1+N)^a$, allowing the constant and the finite exponent to increase
between appearances.  Its coefficients may depend on the fixed
prepared data and on $\epsilon$, but never on the outer truncation
radius.  A bound merely finite for each fixed $N$ will not be denoted
by $\mathcal P_N$.

\begin{proposition}[Coercivity]\label{sy:coercivity}
For all the variables in Equation~\ref{sy:variables},
\begin{equation}\label{sy:coercivity-estimate}
 |T|^2+|M|^2+dj^2+F^2+|dt|_{A_d}^2
 \le20(1+N)^2\mathcal T.
\end{equation}
At a point where $dt=0$, set
$\mathcal E_0=|T^0|^2+|M|^2+F^2+\chi j^2$.
Then, independently of $N$,
\begin{equation}\label{sy:stationary-coercivity}
 \tfrac12\mathcal E_0\le\mathcal T\le\mathcal E_0.
\end{equation}
\end{proposition}

\begin{proof}
Completing first the $x$ square and then the $F$ square in
Equation~\ref{sy:expanded-T} gives the exact expression
\begin{equation}\label{sy:positive-squares}
\begin{aligned}
 \mathcal T={}&\tfrac12|T^0|^2+|M+s_pa_wy|^2
 +(2s_p-v/4)|y|^2\\
 &+2s_pd\left(x+\frac{a_wj}{q}\right)^2
 +\tfrac34\left(F-\frac{\chi j}{3}\right)^2
 +\frac{d(8+v)}{3q^2}j^2.
\end{aligned}
\end{equation}
Indeed, the residual coefficient of $j^2$ is
\[
 \chi-\frac13\chi^2-s_p\frac{\chi^2v}{2d}
 =\frac{\chi}{q}\left(\frac43+\frac v6\right)
 =\frac{d(8+v)}{3q^2}.
\]
Every displayed square has a positive coefficient, and
$2s_p-v/4\ge3s_p/2\ge3/4$.
The last square implies
$dj^2\le3(N+2)^2\mathcal T/8$ and
$\chi^2j^2\le3\mathcal T/2$.
The $F$ square then gives $F^2\le3\mathcal T$; hence
$|T|^2=|T^0|^2+(F-\chi j)^2/2\le13\mathcal T/2$.
The transverse square gives $|y|^2\le4\mathcal T/3$.
Undoing the $x$ shift gives
\[
 dx^2\le2d\left(x+\frac{a_wj}{q}\right)^2
            +\frac{2dv}{q^2}j^2
 \le\frac{8}{3}\mathcal T.
\]
Likewise,
\[
 |M|^2\le2|M+s_pa_wy|^2+2s_p^2v|y|^2
 \le\left(2+\frac43(N+\tfrac12)\right)\mathcal T.
\]
Adding these bounds proves Equation~\ref{sy:coercivity-estimate}.
In particular the coercivity loses no exponential factor from
$l$ or $d$.

When $dt=0$, Equation~\ref{sy:expanded-T} reduces to
\[
 \mathcal T=\tfrac12|T^0|^2+|M|^2
 +\tfrac34F^2-\tfrac12\chi Fj
 +(\chi-\tfrac14\chi^2)j^2.
\]
The matrix of the last two variables $(F,\sqrt\chi j)$ is
\[
 \begin{pmatrix}3/4&-\sqrt\chi/4\\
                 -\sqrt\chi/4&1-\chi/4\end{pmatrix}.
\]
Its eigenvalues are $1$ and $(3-\chi)/4\in[1/2,3/4]$.
This proves Equation~\ref{sy:stationary-coercivity}.
\end{proof}

For a symmetric tensor $B$, define
\begin{equation}\label{sy:S-definition}
 \mathcal S(B)=\tfrac12\left\{
     (\operatorname{tr}_{A_d}B)^2
       -|B|_{A_d\otimes A_d}^2\right\}.
\end{equation}

\begin{lemma}[Identity and subtraction at a stationary point]
\label{sy:stationary-minimum}
At any point where $dt=0$,
\begin{equation}\label{sy:stationary-identity}
\begin{aligned}
 \tfrac12e^{2t}R_{\widehat g}
 ={}&\tfrac12R_g-v\operatorname{Ric}_g(e,e)+\mathcal S(H^f)\\
 &-vp\operatorname{tr}_{e^\perp}\nabla^2t
       -2\operatorname{tr}_{A_d}\nabla^2t,
\end{aligned}
\end{equation}
and
\begin{equation}\label{sy:stationary-subtraction}
 \mathcal T-\mathcal S(H^f)
 \ge\tfrac1{12}\mathcal T
          -(6N+25)\bigl(vF^2+|K|^2\bigr).
\end{equation}
In particular both Hessian terms in
Equation~\ref{sy:stationary-identity} are nonpositive at a local
minimum of $t$.
\end{lemma}

\begin{proof}
Use the graph of $f$ in the Riemannian warped product
$g+l^2dz^2$.  At $dt=0$ its second fundamental form, up to its
irrelevant overall normal sign, is $H^f$; its induced metric is
$\bar g$ and its inverse is $A_d$.
At that point $l^{-1}\nabla^2l=p\nabla^2t$.
The horizontal Ricci tensor of the ambient metric is
$\operatorname{Ric}_g-p\nabla^2t$, its unit vertical Ricci component
is $-p\Delta_gt$, and its mixed Ricci components vanish.  Its scalar
curvature is $R_g-2p\Delta_gt$.
The graph unit normal has horizontal part $-w$ and vertical part of
length $\sqrt d$.  The contracted graph Gauss equation therefore gives
\[
 \tfrac12R_{\bar g}
 =\tfrac12R_g-v\operatorname{Ric}_g(e,e)+\mathcal S(H^f)
      -vp\operatorname{tr}_{e^\perp}\nabla^2t.
\]
At a stationary point the conformal law contributes exactly
$-2\operatorname{tr}_{A_d}\nabla^2t$.
This proves Equation~\ref{sy:stationary-identity}, including all
the Hessian terms.

To prove the inequality, first use $C=H^f+K$.
Its transverse trace is $F-\chi j$, so direct contraction gives
\[
 \mathcal S(C)=\tfrac14(F-\chi j)^2-\tfrac12|T^0|^2
                     +dj(F-\chi j)-d|M|^2.
\]
Subtracting this from the stationary expression for $\mathcal T$
gives the exact formula
\begin{equation}\label{sy:exact-subtraction}
\begin{aligned}
 \mathcal T-\mathcal S(C)
 ={}&|T^0|^2+(1+d)|M|^2+\tfrac12F^2-dFj\\
 &+b_*j^2,\qquad b_*=\chi(1+d)-\tfrac12\chi^2.
\end{aligned}
\end{equation}
Since $\chi\le d$, we have $b_*\ge\chi(1+d/2)\ge\chi$ and
\[
 \frac{d^2}{2b_*}
 \le\frac{dq}{2(2+d)}
 \le\frac13+\frac{pv}{6}.
\]
Young's inequality, in the form
$d|Fj|\le b_*j^2/2+d^2F^2/(2b_*)$, thus proves
\begin{equation}\label{sy:subtraction-before-K}
 \mathcal T-\mathcal S(C)
 \ge\tfrac16\mathcal E_0-\frac p6vF^2.
\end{equation}

For the remaining polarization, let
$X=A_d^{1/2}CA_d^{1/2}$ and
$Y=A_d^{1/2}KA_d^{1/2}$, regarding these as symmetric matrices in a
$g$-orthonormal frame.  Then
\[
 |X|^2=|T^0|^2+\tfrac12(F-\chi j)^2+2d|M|^2+d^2j^2
 \le\frac{N+4}{2}\mathcal E_0,
 \qquad |Y|\le|K|.
\]
Here $d^2/\chi=dq/2\le(N+2)/2$ controls the only axial loss.
The trace-reversal map on symmetric three-by-three matrices has
operator norm two, whence
\[
 |(\operatorname{tr}X)(\operatorname{tr}Y)-\langle X,Y\rangle|
 \le\sqrt{2(N+4)}\sqrt{\mathcal E_0}|K|.
\]
Also $\mathcal S(K)\le|K|^2$.  As
$\mathcal S(C-K)=\mathcal S(C)
 -\{(\operatorname{tr}X)(\operatorname{tr}Y)-\langle X,Y\rangle\}
 +\mathcal S(K)$, Young's inequality with coefficient $1/12$ of
$\mathcal E_0$ and Equation~\ref{sy:subtraction-before-K} imply
\[
 \mathcal T-\mathcal S(H^f)
 \ge\tfrac1{12}\mathcal E_0-\tfrac p6vF^2-(6N+25)|K|^2.
\]
Use $\mathcal E_0\ge\mathcal T$ and $0\le p\le N$ to obtain
Equation~\ref{sy:stationary-subtraction}.
\end{proof}

\subsection{The equations and their continuation family}

Choose
\begin{equation}\label{sy:weights}
 \frac34<b_1<1,\qquad
 \rho_0=r^{-3-\beta},\qquad \rho_1=r^{-2b_1},\qquad
 \rho=c_*\rho_0,
\end{equation}
where $c_*>0$ is sufficiently small that the right side of
Equation~\ref{sy:strict-margin} is strictly larger than $2\rho$ on
the original exterior.  These weights are extended positively over
the filling by the choice of $r$ in Lemma~\ref{sy:filled-setup}.
Choose a fixed $0<\delta_0<1/24$, and put
\begin{equation}\label{sy:penalty}
 m_0(t)=\vartheta(N(t+\epsilon)),\qquad
 \Xi=\delta_0\bigl(\mathcal T+\eta a_w|df|\bigr)+\rho
                 -m_0(t)C_N(\rho_0+v\rho_1).
\end{equation}
The constant $C_N$ is positive and polynomially bounded in $1+N$.
It is chosen large enough for the floor estimate in
Lemma~\ref{es:floor}; that lemma will show that this is a permissible
polynomial choice.  Thus the parameter order is: the prepared
exterior, $\epsilon$, $N$, the above coefficients, and finally a
sufficiently large truncation radius $R$.  At $t=-\epsilon$ the
penalty multiplier is one, while it vanishes at
$t\ge-\epsilon+1/N$.

On $\Omega_{N,R}$, obtained by cutting the sole end at $r=R$, the
system is
\begin{equation}\label{sy:system}
 F=h=\eta f,\qquad
 \operatorname{div}_gV=u\Xi,
 \qquad f=Z=0\quad\hbox{on }\{r=R\}.
\end{equation}
The first equation is equivalently
\begin{equation}\label{sy:normalized-trace}
 \operatorname{tr}_{A_\chi}\nabla^2f
 =\frac{\sqrt D}{l}
       \bigl(\eta f-\operatorname{tr}_{A_\chi}K\bigr).
\end{equation}
The penalty uses the smooth quantity
$\eta a_w|df|=\eta lD^{-1/2}|df|^2$ even at zero slope.
All quantities in Equations~\ref{sy:system} and
\ref{sy:normalized-trace} are now specified; the original boundary
lies in the interior of $\Omega_{N,R}$.

For the continuation argument, let
\begin{equation}\label{sy:homotopy}
 (a,b)\in\mathcal H
 :=([0,1]\times\{1\})\cup(\{0\}\times[0,1]),
 \qquad K_a=aK,
 \qquad F_a=\eta f,
 \qquad \operatorname{div}V_a=b\,u\Xi_a.
\end{equation}
Every expression containing $K$, including $F,V,\mathcal T,\tau$
and the constraint densities in scalar identities, is recomputed
with $K_a$.  The weights, $p,l,\eta$, $\delta_0$, and $C_N$ remain
fixed.  The original system is $(a,b)=(1,1)$; the two segments lead
first to $(0,1)$ and then to $(0,0)$.

\begin{lemma}[The zero-tensor segment]\label{sy:zero-trace-reduction}
Every smooth zero-Dirichlet solution of the trace equation with $K=0$
has $f=0$.  On that segment $t=Z$, $d=\chi=1$, $w=0$, and the
second equation becomes
\begin{equation}\label{sy:scalar-end-system}
 2\Delta_gZ+2(1+p(Z))|dZ|^2
 =b\{2\delta_0s_p(Z)|dZ|^2+\rho-m_0(Z)C_N\rho_0\}.
\end{equation}
In particular, at $(a,b)=(0,0)$ the unique solution is $f=Z=0$.
\end{lemma}

\begin{proof}
At a positive interior maximum of $f$, the left side of
Equation~\ref{sy:normalized-trace} is nonpositive and the right side
is strictly positive.  A negative minimum is excluded in the same
way.  The zero boundary value therefore gives $f=0$.
The defining formulas now give $t=Z$,
$\mathcal T=2s_p|dZ|^2$, $u=e^Zl(Z)$, and
$V=2e^Zl(Z)\nabla Z$.  Differentiating this flux proves
Equation~\ref{sy:scalar-end-system}.
For $b=0$, define
$\Psi(z)=\int_0^z e^sl(s)\,ds$.  The equation is
$\Delta_g\Psi(Z)=0$ with zero boundary value, so
$\Psi(Z)=0$ by the maximum principle.  Since $\Psi'>0$, $Z=0$.
This identifies the starting solution for continuation without
assuming existence at any other parameter value.
\end{proof}
